\section{Paper Metadata and Overview}

\paragraph{Reference.}
Song Mei, Andrea Montanari, and Phan-Minh Nguyen,
\emph{A mean field view of the landscape of two-layer neural networks},
PNAS 115(33):E7665--E7671, 2018 \cite{mei2018meanfield}.

\paragraph{Main problem.}
The paper studies two-layer networks trained by stochastic gradient descent (SGD) in the population setting. The main claim is that, in a suitable large-width / small-step-size regime, the empirical parameter distribution evolves according to a nonlinear PDE. This PDE is a Wasserstein gradient flow, and that description simplifies the loss-landscape analysis.

\paragraph{Core mathematical objects.}
The network prediction is
\begin{equation}
\hat y(x;\theta)=\frac{1}{N}\sum_{i=1}^N \sigma_*(x;\theta_i),
\end{equation}
with population risk
\begin{equation}
R_N(\theta)=\E\bigl[(y-\hat y(x;\theta))^2\bigr].
\end{equation}
After rewriting the risk in terms of the empirical measure
\(
\hat \rho^{(N)} = N^{-1}\sum_{i=1}^N \delta_{\theta_i},
\)
the paper introduces the mean-field functional
\begin{equation}
R(\rho)=R_\# + 2\int V(\theta)\,\rho(\dd \theta)
+ \int U(\theta_1,\theta_2)\,\rho(\dd\theta_1)\rho(\dd\theta_2).
\end{equation}
The corresponding distributional dynamics (DD) is
\begin{equation}
\partial_t \rho_t = 2\xi(t)\,\nabla_\theta \cdot \Bigl(\rho_t \nabla_\theta \Psi(\theta;\rho_t)\Bigr),
\qquad
\Psi(\theta;\rho)=V(\theta)+\int U(\theta,\theta')\,\rho(\dd\theta').
\end{equation}

\paragraph{High-level message.}
The proof spine has three layers:
\begin{enumerate}[label=\arabic*.]
\item rewrite the finite-width problem as an empirical-measure problem;
\item show SGD converges to a PDE by a propagation-of-chaos argument;
\item analyze the PDE rather than the original nonconvex optimization problem.
\end{enumerate}
The PDE can be studied using gradient-flow tools and, in some examples, yields near-global convergence guarantees.

\paragraph{Role of the Supplementary Information.}
The Supplementary Information is not optional here. It contains most of the formal proof architecture: Corollary~2.1, Lemmas in Sections~3--6, the proof of Theorem~3, and the detailed PDE / free-energy arguments that support Theorems~4--7.
